\documentclass[../../main.tex]{subfiles}
\begin{document}

\begin{center}
  {\fontsize{48pt}{58pt}\selectfont \textbf{1961}\normalsize\textbf{年度}}
\end{center}

\begin{table}[htb]
  \centering
  \caption{出題テーマと難易度}
  \begin{tabular}{|c|p{9cm}|c|}
    \hline
    \textbf{問題番号} & \textbf{テーマ} & \textbf{難易度} \\
    \hline
    第1問 & 三角関数と証明 & \\
    \hline
    第2問 & 三角関数の最大最小 & \\
    \hline
    第3問 & 三角関数と最大最小 & \\
    \hline
    第4問 & 軌跡（図形と方程式） & \\
    \hline
    第5問 & 積分と不等式 & \\
    \hline
    第6問 & 関数方程式と証明 & \\
    \hline
  \end{tabular}
\end{table}

\begin{table}[htb]
  \centering
  \caption{解答}
  \renewcommand{\arraystretch}{1.4}
  \begin{tabular}{|c|c|p{9cm}|}
    \hline
    \textbf{問題番号} & \textbf{小問} & \textbf{答え} \\
    \hline
    第1問 & - & 証明問題 \\
    \hline
    第2問 & - & $-\sqrt{2} \le a \le 4+3\sqrt{2}$ \\
    \hline
    第3問 & 最大 & $(x,y)=\left(\dfrac{\sqrt{3}}{2},\dfrac{1}{2}\right),\left(-\dfrac{\sqrt{3}}{2},-\dfrac{1}{2}\right)$ \\
    \cline{2-3}
    & 最小 & $(x,y)=\left(-\dfrac{1}{2},\dfrac{\sqrt{3}}{2}\right),\left(\dfrac{1}{2},-\dfrac{\sqrt{3}}{2}\right)$ \\
    \hline
    第4問 & - & $X^2+4Y^2=1$ \\
    \hline
    第5問 & - & $p^2 \le 3q$ \\
    \hline
    第6問 & (1) & 証明問題 \\
    \cline{2-3}
    & (2) & 証明問題 \\
    \hline
  \end{tabular}
\end{table}

\paragraph{試験形態}
6問．試験時間は要確認．

\end{document}
